\chapter{Perbandingan Fungsi}\label{ch:b2:comparison}

Analisis asimtotik --- seni mengganti besaran yang rumit
dengan besaran sederhana ditambah galat yang terkendali --- dimulai pada
jilid Tahun ke-1 lewat uraian Taylor. Bab ini menjadikannya disiplin
tersendiri: uraian sepanjang skala umum, perbandingan deret dengan
integral beserta seluruh daya asimtotiknya, rumus Stirling (yang
dibuktikan tuntas), dan telaah sistematis atas barisan yang terdefinisi
secara implisit. Teknik itu adalah santapan sehari-hari analisis
asimtotik, dan setiap bab berikutnya yang menaksir apa pun ---
deret, integral, peluang --- makan dari meja ini.

\section{Relasi perbandingan dan skala}

\begin{definition}\label{def:b2:comparison:landau}
Di dekat sebuah titik $a$ ($a \in \R$ atau $\pm\infty$), untuk fungsi
(atau barisan, dengan $n \to \infty$): $f = o(g)$, $f = O(g)$, $f \sim g$
seperti pada jilid Tahun ke-1. Sebuah \emph{skala
perbandingan}\index{skala perbandingan} di $a$ adalah keluarga fungsi
positif yang sebanding sepasang demi sepasang dan terurut total oleh
$o(\cdot)$ --- skala bakunya di $+\infty$ adalah
\[
x^{\alpha} (\ln x)^{\beta}
\qquad (\alpha, \beta \in \R),
\]
yang terurut secara leksikografis pada $(\alpha, \beta)$, dan diperhalus
bila perlu oleh eksponensial $\eu^{\gamma x}$.
\end{definition}

\begin{definition}[Uraian asimtotik]\label{def:b2:comparison:expansion}
Fungsi $f$ menerima \emph{uraian asimtotik}\index{uraian
asimtotik}
\[
f = c_1 \varphi_1 + c_2\varphi_2 + \dots + c_k \varphi_k +
o(\varphi_k)
\qquad (\varphi_{i+1} = o(\varphi_i) \text{ pada skalanya})
\]
apabila sisa berturut-turutnya memenuhi taksiran di atas.
Koefisiennya lalu tunggal: $c_1 = \lim f/\varphi_1$, dan
secara induktif $c_{i+1} = \lim\,(f - \sum_{j \leq i}
c_j\varphi_j)/\varphi_{i+1}$.
\end{definition}

\begin{example}
Uraian Taylor adalah \omterm{def:b2:comparison:expansion}{uraian asimtotik} sepanjang skala $(x -
a)^k$ di $a$. Namun gagasannya jauh lebih luas: di $+\infty$,
\[
\frac{1}{x - \ln x}
= \frac1x \cdot \frac{1}{1 - \frac{\ln x}{x}}
= \frac1x + \frac{\ln x}{x^2} + o\Bigl(\frac{\ln x}{x^2}\Bigr),
\]
yaitu uraian sepanjang skala campuran --- tak ada teorema Taylor yang
berlaku, hanya uraian geometri dan kalkulus atas $o$.
\end{example}

\begin{example}[Skala bakunya sungguh terurut]\label{ex:b2:comparison:scaleorder}
Klaim leksikografis pada \cref{def:b2:comparison:landau}
memerlukan satu baris bukti untuk tiap kasusnya. Bandingkan
$x^{\alpha}(\ln x)^{\beta}$ dan $x^{\alpha'}(\ln x)^{\beta'}$
di $+\infty$. Jika $\alpha < \alpha'$: rasionya adalah
$x^{\alpha - \alpha'}(\ln x)^{\beta - \beta'} \to 0$, sebab
pangkat negatif $x$ meremukkan pangkat berapa pun dari $\ln x$ (ambil $x =
\eu^t$: $\eu^{(\alpha - \alpha')t}\,t^{\beta - \beta'} \to 0$
lewat limit ``eksponensial mengalahkan polinomial'' pada jilid Tahun
ke-1). Jika $\alpha = \alpha'$ dan $\beta < \beta'$: rasionya
langsung $(\ln x)^{\beta - \beta'} \to 0$. Jadi pasangan
$(\alpha, \beta)$, yang terurut secara leksikografis, mengurutkan
skalanya oleh $o(\cdot)$ --- dan penyulihan $x = \eu^t$ adalah
siasat serbaguna untuk perbandingan campuran pangkat dan logaritma.
\end{example}

\begin{example}[Mengurutkan satu kebun binatang]\label{ex:b2:comparison:ranking}
Skala wajib \emph{terurut}; berikut latihan bakunya. Di
$+\infty$, bandingkan $n^{10}$, $\eu^{\sqrt{\ln n}\,\cdot\,\sqrt n}$,
$2^n$ dan $n^{\ln n}$ dengan mengambil logaritmanya:
\[
10\ln n
\;\ll\; (\ln n)^2
\;\ll\; \sqrt{n\ln n}
\;\ll\; n\ln 2 ,
\]
dengan $a_n \ll b_n$ berarti $a_n = o(b_n)$; adapun entri keduanya
adalah $\ln(n^{\ln n})$. Eksponensial memelihara jurang tegas itu (sebab
jika $\ln u_n - \ln v_n \to -\infty$ maka $u_n/v_n \to 0$), sehingga
\[
n^{10} = o\bigl(n^{\ln n}\bigr),
\qquad
n^{\ln n} = o\bigl(\eu^{\sqrt{n\ln n}}\bigr),
\qquad
\eu^{\sqrt{n\ln n}} = o(2^n) .
\]
Pelajarannya ada dua: selalulah membandingkan lewat logaritma
(yakni selisih logaritma, bukan rasio logaritma), dan jangan pernah
menyimpulkan $u_n \sim v_n$ dari $\ln u_n \sim \ln v_n$ --- sebab pasangan
$n^{10}$ dan $n^{\ln n}$ punya rasio $\ln$ yang menuju $\infty$, sedangkan
$2^n$ dan $4^n$ punya rasio $\ln$ tepat $2$ padahal keduanya sama sekali
tidak setara.
\end{example}

\section{Perbandingan deret dengan integral, secara asimtotik}

\begin{theorem}\label{thm:b2:comparison:seriesintegral}
Misalkan $f$ \omterm{def:b2:metric:continuity}{kontinu}, positif, dan \emph{turun} pada
$\intco{1}{+\infty}$.
\begin{enumerate}
  \item Jika $\int_1^{\infty} f$ konvergen, maka sisanya memenuhi
        \[
        \int_{n+1}^{\infty} f \;\leq\; \sum_{k > n} f(k) \;\leq\;
        \int_{n}^{\infty} f .
        \]
  \item Jika $\int_1^\infty f$ divergen, maka jumlah parsialnya memenuhi
        $\sum_{k=1}^{n} f(k) = \int_1^n f + C + o(1)$ untuk suatu
        konstanta $C$: yakni selisih $\sum_{k \leq n} f(k) -
        \int_1^n f$ \emph{konvergen}.
\end{enumerate}
\end{theorem}

\begin{proof}
Pengapitan $f(k+1) \leq \int_k^{k+1} f \leq f(k)$ (dari keturunannya) tak
lain alat Tahun ke-1; menjumlahkannya atas $k \geq n+1$, atau atas $k
\geq n$, memberi (1). Untuk (2), tetapkan $u_k = f(k) - \int_k^{k+1} f$:
menurut pengapitan itu, $0 \leq u_k \leq f(k) - f(k+1)$, jadi jumlah
parsial $\sum u_k$ terbatas oleh teleskop $f(1) - f(n+1) \leq
f(1)$: sehingga deretnya konvergen. Lebih lanjut barisan
$\bigl(\int_n^{n+1} f\bigr)_n$ tidak naik (sebab $f$ turun) dan tak
negatif, jadi ia konvergen. Dengan menulis
\[
\sum_{k=1}^{n} f(k) - \int_1^n f
= \sum_{k=1}^{n} u_k + \int_n^{n+1} f ,
\]
ruas kanannya konvergen ketika $n \to \infty$: jadi selisihnya konvergen
ke sebuah konstanta $C$, dan itulah pernyataan (2).
\end{proof}

\begin{example}[Uraian harmonik]\label{ex:b2:comparison:harmonic}
Untuk $f(t) = \frac1t$: $H_n = \ln n + \gamma + o(1)$, yang memulihkan
konstanta Euler (jilid Tahun ke-1) dengan bukti yang lebih bersih.
Didorong satu orde lebih jauh (\cref{exo:b2:comparison:3}):
\[
H_n = \ln n + \gamma + \frac{1}{2n} + o\Bigl(\frac1n\Bigr).
\]
Angkanya membuat keuntungannya kasatmata di $n = 10$: $H_{10} =
2.928968\dots$ dan $\ln 10 = 2.302585\dots$, jadi taksiran mentah
$\gamma$ adalah $H_{10} - \ln 10 = 0.626383$, yang meleset $0.049$;
setelah dikurangi koreksi $\frac1{20}$ diperoleh $0.576383$,
yang meleset dari $\gamma = 0.577216$ hanya $8.3\cdot10^{-4}$ ---
dan angka itu sendiri adalah suku berikutnya $\frac{1}{12\cdot100}$ pada
uraiannya, sebagaimana dibuktikan soal akhir pekan (pertanyaan 8).
\end{example}

\begin{example}[Taksiran kasar $\ln(n!)$ tanpa Stirling]\label{ex:b2:comparison:lnfactorial}
Alat pengapitan itu saja sudah menempatkan $\ln(n!)$. Karena
$\ln$ naik,
\[
\int_{k-1}^{k}\ln t\,\dd t \;\leq\; \ln k \;\leq\;
\int_{k}^{k+1}\ln t\,\dd t ,
\]
lalu menjumlahkannya atas $k = 2, \dots, n$ (dengan $\int_1^n\ln = n\ln n
- n + 1$) memberi
\[
n\ln n - n + 1 \;\leq\; \ln(n!) \;\leq\; (n+1)\ln(n+1) - n .
\]
Kedua pagarnya berbentuk $n\ln n - n + O(\ln n)$: karenanya $\ln(n!) =
n\ln n - n + O(\ln n)$, dan khususnya $\ln(n!) \sim n\ln n$.
Yang ditambahkan Stirling adalah dua anak tangga berikutnya --- yaitu
$\frac12\ln n$ dan konstanta $\ln\sqrt{2\pi}$ --- yang berharga
teleskop lebih halus pada \cref{thm:b2:comparison:stirling}. Mengetahui
ketelitian yang dibeli tiap perkakas adalah separuh dari seni asimtotik.
\end{example}

\begin{example}[Pencoretan menuntut uraian]\label{ex:b2:comparison:cancellation}
Hitunglah limit $\sqrt{n^2 + n} - n$. Kedua sukunya
$\sim n$, dan ``$\sim n - n$'' tak bermakna: sebab yang setara
tak dapat dikurangkan. Uraikanlah sebagai gantinya:
\[
\sqrt{n^2 + n} - n
= n\Bigl(\sqrt{1 + \tfrac1n} - 1\Bigr)
= n\Bigl(\frac{1}{2n} - \frac{1}{8n^2} +
O\Bigl(\frac1{n^3}\Bigr)\Bigr)
= \frac12 - \frac{1}{8n} + O\Bigl(\frac1{n^2}\Bigr) :
\]
jadi limitnya $\frac12$, dengan laju penghampiran $\frac1{8n}$ sebagai
bonus. Mekanismenya pantas dinamai: selisih dua besaran besar
yang setara hidup seluruhnya pada suku \emph{berikutnya},
jadi kita wajib menguraikannya sampai orde pertama tempat kedua
ruasnya berselisih --- lalu membawa sisanya untuk mengesahkan bahwa
tak ada lagi yang bertahan pada orde itu.
\end{example}

\begin{example}[Sebuah perbandingan divergen, dikerjakan]\label{ex:b2:comparison:loglog}
Untuk $f(t) = \frac{1}{t\ln t}$ pada $\intco{2}{+\infty}$
(yang \omterm{def:b2:metric:continuity}{kontinu}, positif, turun): $\int_2^x f = \ln\ln x -
\ln\ln 2 \to \infty$, jadi menurut
\cref{thm:b2:comparison:seriesintegral}~(2),
\[
\sum_{k=2}^{n}\frac{1}{k\ln k} = \ln\ln n + C + o(1)
\]
untuk suatu konstanta $C$. Ada dua pelajaran. Pertama, kedivergenannya
nyata tetapi selambat gletser: jumlah parsialnya baru melampaui $4$
sekitar $n \approx \eu^{\eu^{4 - C}}$, yang besarnya astronomis. Kedua,
\emph{bentuk} $\ln\ln n$ itu diserahkan oleh sebuah antiturunan, bukan
ditebak: sebab untuk suku yang monoton, integral adalah alat penjumlah
yang kanonik, dan konstanta $C$ --- seperti $\gamma$ milik Euler
--- adalah ingatan atas suku-suku awalnya.
\end{example}

\section{Rumus Stirling}

\begin{lemma}[Integral Wallis, ditinjau ulang]\label{lem:b2:comparison:wallis}
Misalkan $W_n = \int_0^{\pi/2} \sin^n t\,\dd t$. Maka $nW_nW_{n-1} =
\frac\pi2$ untuk $n \geq 1$, barisan $(W_n)$ turun, dan $W_n \sim
\sqrt{\dfrac{\pi}{2n}}$.
\end{lemma}

\begin{proof}
Pengintegralan parsial memberi $nW_n = (n-1)W_{n-2}$ ($n \geq 2$), jadi
$nW_nW_{n-1}$ tetap terhadap $n$, sama dengan $1 \cdot W_1 W_0 =
\frac\pi2$. Keturunannya: $\sin^{n+1} \leq \sin^n$ pada
$\intcc{0}{\frac\pi2}$. Pengapitannya, secara rinci: kemonotonan
memberi $W_{n+1} \leq W_n \leq W_{n-1}$, lalu setelah dibagi $W_{n-1}
> 0$,
\[
\frac{n}{n+1} = \frac{W_{n+1}}{W_{n-1}} \leq
\frac{W_n}{W_{n-1}} \leq 1 ,
\]
dengan kesamaan di kirinya berasal dari rekurensi pada indeks $n + 1$.
Kedua batasnya menuju $1$: jadi $W_n \sim W_{n-1}$, sehingga
\[
nW_n^2 \sim nW_nW_{n-1} = \frac\pi2
\qquad\Longrightarrow\qquad
W_n \sim \sqrt{\frac{\pi}{2n}} .
\]
\end{proof}

\begin{example}[Integral Wallis yang pertama]\label{ex:b2:comparison:wallisvalues}
Dari $W_0 = \frac\pi2$, $W_1 = 1$ dan rekurensi $nW_n =
(n-1)W_{n-2}$:
\[
W_2 = \frac\pi4, \qquad
W_3 = \frac23, \qquad
W_4 = \frac{3\pi}{16}, \qquad
W_5 = \frac{8}{15}, \qquad
W_6 = \frac{5\pi}{32}.
\]
Indeks genapnya mengusung $\pi$, indeks ganjilnya rasional --- yakni dua
hasil kali berselang-seling pada bentuk tertutupnya. Secara numerik $W_6
\approx 0.4909$ terhadap nilai asimtotiknya
$\sqrt{\pi/12} \approx 0.5116$: jadi pada $n = 6$ yang setara itu
sudah berselisih di bawah $5\%$, sedangkan kesamaan hasil kalinya eksak
di setiap $n$: $6\,W_6W_5 = 6\cdot\frac{5\pi}{32}\cdot\frac8{15} =
\frac\pi2$. Tabel kecil seperti ini adalah cara termurah untuk
menangkap keseleo aljabar sebelum ia sempat menulari suatu hujah asimtotik.
\end{example}

\begin{theorem}[Stirling]\label{thm:b2:comparison:stirling}
\[
n! \;\sim\; \sqrt{2\pi n}\, \Bigl(\frac{n}{\eu}\Bigr)^{\!n}
.\index{rumus Stirling}
\]
\end{theorem}

\begin{proof}
\emph{Langkah 1: $n! \sim C \sqrt n\, (n/\eu)^n$ untuk suatu konstanta $C
> 0$.} Tetapkan
\[
d_n = \ln(n!) - \Bigl(n + \frac12\Bigr)\ln n + n .
\]
Maka
\[
d_n - d_{n+1}
= \Bigl(n + \frac12\Bigr) \ln\frac{n+1}{n} - 1
= \Bigl(n + \frac12\Bigr)\Bigl(\frac1n - \frac{1}{2n^2} +
\frac{1}{3n^3} + o\bigl(n^{-3}\bigr)\Bigr) - 1
= \frac{1}{12n^2} + o\Bigl(\frac{1}{n^2}\Bigr),
\]
lewat uraian Taylor $\ln(1 + \frac1n)$. Karenanya deret $\sum (d_n
- d_{n+1})$ konvergen mutlak (dengan pembandingan terhadap $\sum
n^{-2}$), sehingga $(d_n)$ konvergen, katakanlah ke $d$; setelah
dieksponensialkan, $n! \sim C\sqrt n\,(n/\eu)^n$ dengan $C = \eu^{d}$.

\emph{Langkah 2: $C = \sqrt{2\pi}$ lewat Wallis.} Bentuk tertutup
$W_{2p} = \frac{(2p)!}{4^p (p!)^2}\cdot\frac\pi2$ (dari rekurensinya,
yaitu perhitungan Tahun ke-1 yang dikerjakan ulang pada latar
\cref{lem:b2:comparison:wallis}) berpadu dengan Langkah 1:
\[
W_{2p} \sim \frac{C\sqrt{2p}\,(2p/\eu)^{2p}}
{4^p\,\bigl(C\sqrt p\,(p/\eu)^p\bigr)^2}\cdot\frac{\pi}{2}
= \frac{\sqrt{2p}}{C\,p}\cdot\frac{\pi}{2}
= \frac{\pi}{C}\cdot\frac{1}{\sqrt{2p}} .
\]
Dibandingkan dengan $W_{2p} \sim \sqrt{\frac{\pi}{4p}}$
(\cref{lem:b2:comparison:wallis}): syarat $\frac{\pi}{C\sqrt{2p}} =
\sqrt{\frac{\pi}{4p}}\,(1 + o(1))$ memaksa $C = \pi
\sqrt{\frac{4p}{2p\,\pi}} = \sqrt{2\pi}$.
\end{proof}

\begin{example}[Koefisien binomial pusat]\label{ex:b2:comparison:centralbinomial}
\[
\binom{2n}{n} = \frac{(2n)!}{(n!)^2}
\sim \frac{\sqrt{4\pi n}\,(2n/\eu)^{2n}}{2\pi n\,(n/\eu)^{2n}}
= \frac{4^n}{\sqrt{\pi n}} :
\]
peluang bahwa jalan acak yang setangkup kembali ke $0$ pada saat
$2n$ bernilai $\sim \frac{1}{\sqrt{\pi n}}$ --- yakni pengumuman awal
bagi \cref{ch:b2:randomvar}.
\end{example}

\begin{remark}[Pandangan ke depan di dalam jilid ini]\label{rem:b2:comparison:perspectives}
Setiap bab kuantitatif berikutnya berbicara dalam bahasa bab
ini. \cref{ch:b2:series} menggolongkan deret dengan membandingkan
sukunya terhadap skala $n^{-\alpha}(\ln n)^{-\beta}$ --- dan soal
akhir pekannya memetakan perbatasan itu seluruhnya.
\cref{ch:b2:integration} melakukan hal yang sama untuk integral tak
wajar, dengan skala yang persis sama pada peubah \omterm{def:b2:metric:continuity}{kontinunya}.
\cref{ch:b2:powerseries} menghitung jari-jari kekonvergenan dari
$\limsup\abs{a_n}^{1/n}$, yakni latihan mencari yang setara bagi akar
ke-$n$, dan Stirling menjadi kunci bakunya ($\sqrt[n]{n!}
\sim \frac n\eu$, \cref{exo:b2:comparison:4}). Sedangkan bab
peluangnya menguangkan Stirling secara langsung: taksiran
lokal pada \cref{ch:b2:randomvar} bagi koefisien binomial
tak lain \cref{ex:b2:comparison:centralbinomial} dan
\cref{ex:b2:comparison:lopsided} kata demi kata. Asimtotik bukanlah
sebuah bab di sini; ia logat jilid ini.
\end{remark}

\begin{method}[Daftar periksa bootstrap]\label{met:b2:comparison:checklist}
Sebelum memercayai uraian hasil bootstrap, periksalah empat butir.
(1) \emph{Keberadaan lebih dulu:} akar atau barisannya wajib
terkunci (lewat kemonotonan atau nilai antara) sebelum uraian apa pun
--- sebab lambang tanpa acuan terurai dengan indah dan tak bermakna
apa-apa. (2) \emph{Satu orde tiap lintasan:} tiap penyulihan hanya
boleh dipercaya sampai orde taksiran yang dimasukkan; menarik
dua suku baru dari satu lintasan adalah sumber klasik koefisien yang
salah. (3) \emph{Sisanya ikut menumpang:} bawalah
$o(\cdot)$ melewati setiap langkah aljabarnya lalu biarkan penyerapan
(suku kecil ditelan sisa yang lebih besar) terjadi di
akhir, secara tersurat. (4) \emph{Audit numerik:} nilailah pada satu
nilai $n$ yang jujur; sebab galat koefisien sering sekali selamat dari
penurunan ulang secara aljabar, dan hampir tak pernah selamat dari
aritmetika.
\end{method}

\begin{remark}[Jebakan yang sering muncul]\label{rem:b2:comparison:pitfalls}
(i) Yang setara \emph{menjumlah} dengan buruk: dari $u_n \sim n + \ln n$
dan $v_n \sim -n$ kita \emph{tidak} boleh menyimpulkan $u_n + v_n \sim
\ln n$; sebab pencoretan menuntut uraian dengan sisa yang tersurat,
tak pernah sekadar yang setara. (ii) Jangan pernah mengeksponensialkan
sebuah kesetaraan: $n + 1 \sim n$ tetapi $\eu^{n+1} \not\sim \eu^n$;
arah yang aman adalah mengambil logaritma dari yang setara dan menuju
$+\infty$ (soal akhir pekan bab ini, pertanyaan 24). (iii)
\omterm{def:b2:comparison:expansion}{Uraian asimtotik} melekat pada sebuah \emph{skala}: menulis
$f = \frac1x + o\bigl(\frac1{x^2}\bigr)$ mengklaim lebih banyak daripada
$f = \frac1x + o\bigl(\frac1x\bigr)$, dan mencampur keduanya membatalkan
aljabar sesudahnya. (iv) Dalam bootstrap, sulihkanlah
uraian \emph{seluruhnya} beserta sisanya --- sebab membuang
sebuah $o(\cdot)$ di tengah lintasan menghasilkan koefisien yang masuk
akal tetapi salah. (v) Perbandingan deret dengan integral memerlukan
kemonotonan: untuk suku yang berayun ia gagal sama sekali (bandingkan
$\sum\frac{\sin k}k$, \cref{ch:b2:series}).
\end{remark}

\begin{example}[Stirling dalam angka]\label{ex:b2:comparison:stirlingnum}
Pada $n = 10$: rumusnya memberi $\sqrt{20\pi}\,(10/\eu)^{10}
\approx 3\,598\,696$ terhadap $10! = 3\,628\,800$: jadi galat nisbinya
$8.3\cdot10^{-3}$, sungguh menakjubkan bagi pernyataan ``asimtotik'' di
$n = 10$. Galatnya punya struktur --- yakni perhalusan eksak $n!
= \sqrt{2\pi n}\,(n/\eu)^n\bigl(1 + \frac1{12n} +
O(n^{-2})\bigr)$ --- yang koreksi pertamanya $\frac1{120} \approx
8.3\cdot10^{-3}$ menjelaskan jurang teramati itu hampir persis. Adapun
mesin Euler--Maclaurin pada soal akhir pekan justru sumber sistematis
bagi suku koreksi semacam itu.
\end{example}

\begin{remark}[Di mana bab ini dipakai]
Perbandingan asimtotik adalah tata bahasa segala yang kuantitatif
di hilir: uji kekonvergenan dan panorama Bertrand pada
\cref{ch:b2:series}, kriteria keterintegralan pada
\cref{ch:b2:integration}, perhitungan jari-jari kekonvergenan
pada \cref{ch:b2:powerseries}, dan teorema limit pada
\cref{ch:b2:randomvar} (tempat Stirling menjalankan taksiran de
Moivre--Laplace). Jilid Tahun ke-3 mengindustrikan satu gagasan yang
kita buktikan dengan tangan di sini --- sarikan suku utamanya, batasi
sisanya --- menjadi metode Laplace dan kekonvergenan terdominasi.
\end{remark}

\begin{example}[Sebuah integral dibandingkan dengan dirinya: $\int_2^x \frac{\dd t}{\ln t}$]\label{ex:b2:comparison:liworked}
Kotak perkakas perbandingan ini juga berjalan pada integral. Misalkan
$F(x) = \int_2^x\frac{\dd t}{\ln t}$ (integrannya \omterm{def:b2:metric:continuity}{kontinu} pada
$\intco2\infty$). Integralkan secara parsial:
\[
F(x) = \Bigl[\frac{t}{\ln t}\Bigr]_2^x +
\int_2^x\frac{\dd t}{(\ln t)^2}
= \frac{x}{\ln x} + O\Bigl(\int_2^x\frac{\dd t}{(\ln
t)^2}\Bigr) + O(1),
\]
dan integral sisanya bernilai $o\bigl(\frac{x}{\ln x}\bigr)$:
pecahlah di $\sqrt x$, lalu batasi dengan
\[
\int_2^{\sqrt x}\frac{\dd t}{(\ln t)^2} \leq \sqrt x
\qquad\text{dan}\qquad
\int_{\sqrt x}^{x}\frac{\dd t}{(\ln t)^2} \leq
\frac{x}{(\ln\sqrt x)^2} = \frac{4x}{(\ln x)^2} .
\]
Karenanya $F(x) \sim \frac{x}{\ln x}$. Pembaca yang menjumpai teorema
bilangan prima pada soal akhir pekan bab ini akan mengenali
$F$: ia integral logaritmik, penaksir $\pi(x)$ yang lebih baik, dan
perhitungan tadi menunjukkan ia sepadan dengan
$\frac{x}{\ln x}$ pada orde pertamanya.
\end{example}

\begin{example}[Stirling pada binomial yang berat sebelah]\label{ex:b2:comparison:lopsided}
Rutin tiga faktorial yang sama seperti pada
\cref{ex:b2:comparison:centralbinomial} memberi, untuk
$\binom{3n}{n} = \frac{(3n)!}{n!\,(2n)!}$:
\[
\binom{3n}{n} \sim
\frac{\sqrt{6\pi n}\,(3n/\eu)^{3n}}
{\sqrt{2\pi n}\,(n/\eu)^{n}\cdot\sqrt{4\pi n}\,(2n/\eu)^{2n}}
= \sqrt{\frac{3}{4\pi n}}\,
\Bigl(\frac{27}{4}\Bigr)^{\!n} .
\]
Laju eksponensialnya $\frac{27}4 = \frac{3^3}{2^2}$ tak lain
$\eu^{3n\,H(1/3)}$ dalam notasi entropi pada teori informasi:
jadi binomial yang berat sebelah tumbuh tegas lebih lambat daripada
$4^n$ milik yang di pusat, per dua langkah --- di sini $(27/4)^{1/3}
\approx 1.89 < 2$ per langkah. Setiap asimtotik binomial dalam
kombinatorika dan peluang (\cref{ch:b2:randomvar}) tak lain perhitungan
yang satu ini dengan bobot yang berbeda.
\end{example}

\section{Barisan yang terdefinisi secara implisit}

\begin{method}\label{met:b2:comparison:implicit}
Untuk mencari asimtotik penyelesaian $x_n$ pada persamaan $F(x, n) =
0$:\index{barisan implisit}
\begin{enumerate}
  \item \emph{Setempatkan}: buktikan keberadaan dan ketunggalan $x_n$
        pada interval yang tertentu (lewat kemonotonan atau teorema
        nilai antara), lalu carilah perilaku kasarnya (limit, orde
        pertumbuhannya).
  \item \emph{Bootstrap}: sulihkan bentuk kasarnya $x_n = (\text{suku
        utama})(1 + \varepsilon_n)$ ke dalam persamaannya lalu
        selesaikan untuk orde berikutnya pada $\varepsilon_n$; ulangi,
        dengan tiap lintasan memperhalus satu orde.
\end{enumerate}
\end{method}

\begin{example}\label{ex:b2:comparison:tan}
Untuk $n \geq 1$, persamaan $\tan x = x$ punya tepat satu penyelesaian
$x_n$ di $\intoo{n\pi - \frac\pi2}{n\pi + \frac\pi2}$ (sebab fungsi
$\tan x - x$ naik dari $-\infty$ ke $+\infty$ di sana, dengan
turunannya $\tan^2 x \geq 0$). \emph{Kasarnya:} $x_n = n\pi +
\frac\pi2 - y_n$ dengan $y_n \in \intoo{0}{\pi}$; karena $x_n \to
\infty$ dan $\tan x_n = x_n \to +\infty$, maka $x_n$ mendekati
asimtotnya dari kiri: jadi $y_n \to 0$. \emph{Bootstrap:} $\tan x_n =
\cot y_n = \frac{1}{\tan y_n} \sim \frac{1}{y_n}$, dan persamaan
$\cot y_n = x_n \sim n\pi$ memberi $y_n \sim \frac{1}{n\pi}$. Karenanya
\[
x_n = n\pi + \frac\pi2 - \frac{1}{n\pi} +
o\Bigl(\frac1n\Bigr),
\]
dan prosesnya berlanjut sampai orde berapa pun
(\cref{exo:b2:comparison:6}).
\end{example}

\begin{example}[Metodenya dijalankan kedua kali]\label{ex:b2:comparison:xplusln}
Selesaikan $x + \ln x = n$ secara asimtotik. \emph{Setempatkan:} $x
\mapsto x + \ln x$ naik dari $-\infty$ ke $+\infty$ pada
$\intoo{0}{+\infty}$: jadi ada akar tunggal $x_n$, dan $x_n \to \infty$.
\emph{Kasarnya:} $\ln x_n = o(x_n)$ memberi $x_n \sim n$.
\emph{Bootstrap:} dari $x_n = n - \ln x_n$ dan $\ln x_n = \ln n
+ o(1)$ (logaritma dari yang setara, sebab kedua ruasnya $\to \infty$):
\[
x_n = n - \ln n + o(1) ;
\]
lalu satu lintasan lagi, dengan $\ln x_n = \ln\bigl(n - \ln n + o(1)\bigr) =
\ln n - \frac{\ln n}{n} + o\bigl(\frac{\ln n}n\bigr)$:
\[
x_n = n - \ln n + \frac{\ln n}{n} +
o\Bigl(\frac{\ln n}{n}\Bigr).
\]
(Periksa di $n = 100$: akarnya $x \approx 95.4415$; rumus tiga
sukunya memberi $100 - 4.6052 + 0.0461 = 95.4409$, sedangkan
rumus dua sukunya $95.3948$ --- jadi tiap lintasan meraih orde yang
diramalkan.) Loop yang sama, lanskap ketiga: metode pada
\cref{met:b2:comparison:implicit} tidak peduli seperti apa rupa
persamaannya, asalkan tiap lintasan memisahkan peubah takdiketahui yang
dominan.
\end{example}

\section{Latihan}

\begin{exercise}[$\star$]\label{exo:b2:comparison:1}
Uraikan di $+\infty$, dua suku melampaui suku utamanya:
\[
\sqrt{x^2 + x + 1} ,
\qquad
\ln(x^2 + x) - 2\ln x,
\qquad
\frac{x + \sin x}{x - \ln x} .
\]
\end{exercise}

\begin{exercise}[$\star$]\label{exo:b2:comparison:2}
Berikan sifatnya (konvergen atau divergen) dan, bila divergen,
asimtotik utamanya bagi $\sum_{k \leq n} k^\alpha$ untuk $\alpha >
-1$, $\alpha = -1$, dan $\alpha < -1$, lewat
\cref{thm:b2:comparison:seriesintegral}.
\end{exercise}

\begin{exercise}[$\star\star$]\label{exo:b2:comparison:3}
Buktikan $H_n = \ln n + \gamma + \frac{1}{2n} + o\bigl(\frac1n\bigr)$.
\emph{(Telaahlah $v_n = H_n - \ln n - \gamma$: tunjukkan $v_n - v_{n+1} =
\frac{1}{2n^2} + O(n^{-3})$ lalu jumlahkan ekornya, dengan membandingkan
terhadap $\sum_{k \geq n} \frac{1}{2k^2} \sim \frac{1}{2n}$ ---
\cref{thm:b2:comparison:seriesintegral}~(1).)}
\end{exercise}

\begin{exercise}[$\star\star$]\label{exo:b2:comparison:4}
Dengan Stirling, carilah yang setara bagi: $\dfrac{(3n)!}{(n!)^3}$;
$\;\dfrac{n!}{n^n}$; dan $\;\sqrt[n]{n!}$ (sebagai $\frac n\eu(1 + o(1))$,
yang dicermatkan sampai dua suku).
\end{exercise}

\begin{exercise}[$\star\star$]\label{exo:b2:comparison:5}
Untuk $n \geq 2$, buktikan bahwa $x^n + x = 1$ punya penyelesaian
tunggal $x_n \in \intoo{0}{1}$, bahwa $x_n \to 1$, lalu tegakkan
\[
x_n = 1 - \frac{\ln n}{n} + o\Bigl(\frac{\ln n}{n}\Bigr).
\]
\emph{(Dari $x_n^n = 1 - x_n$: ambil logaritmanya lalu bootstrap dengan
$x_n = 1 - \varepsilon_n$.)}
\end{exercise}

\begin{exercise}[$\star\star$]\label{exo:b2:comparison:6}
Doronglah \cref{ex:b2:comparison:tan} satu orde lebih jauh:
\[
x_n = n\pi + \frac\pi2 - \frac{1}{n\pi} + \frac{1}{2n^2\pi} +
o\Bigl(\frac{1}{n^2}\Bigr).
\]
\emph{(Tulis $\cot y_n = x_n$ secara eksak, uraikan $\cot y = \frac1y -
\frac y3 + o(y)$ dan $x_n = n\pi(1 + \frac{1}{2n} - \dots)$, lalu
samakan.)}
\end{exercise}

\begin{exercise}[$\star\star$]\label{exo:b2:comparison:7}
Tentukan $\lim_{n\to\infty} \dfrac{1}{n!}\sum_{k=0}^{n} k!$
\emph{(batasi jumlah semua sukunya kecuali dua yang terakhir)}, lalu
turunkan \omterm{def:b2:comparison:expansion}{uraian asimtotik} $\sum_{k \leq n} k! = n!\bigl(1 + \frac1n +
O(n^{-2})\bigr)$.
\end{exercise}

\begin{exercise}[$\star\star\star$]\label{exo:b2:comparison:8}
Misalkan $u_0 > 0$ dan $u_{n+1} = u_n + \dfrac{1}{u_n}$. Buktikan bahwa
$u_n \to \infty$, lalu bahwa $u_n \sim \sqrt{2n}$ \emph{(telaahlah
$u_n^2$: pertambahannya $2 + u_n^{-2}$; lalu jumlahkan)}, lalu perhalus:
\[
u_n = \sqrt{2n}\Bigl(1 + \frac{\ln n}{8n} +
o\Bigl(\frac{\ln n}{n}\Bigr)\Bigr).
\]
\emph{(Dari $u_n^2 = 2n + \sum_{k<n} u_k^{-2} + u_0^2$ dan $u_k^2
\sim 2k$: jumlahnya $\sim \frac12\ln n$ menurut
\cref{thm:b2:comparison:seriesintegral}.)}
\end{exercise}

\begin{exercise}[$\star\star\star$]\label{exo:b2:comparison:9}
(Jumlah Riemann dengan sebuah pelintir) Tentukan perilaku asimtotik
\[
S_n = \sum_{k=1}^{n} \frac{1}{n + k\ln n} .
\]
\emph{(Keluarkan faktor $n$: $S_n = \frac1n\sum_k \bigl(1 +
\frac{k\ln n}{n}\bigr)^{-1}$; kenalilah jumlah bergaya Riemann dengan
parameter $t = \ln n$ yang berubah lambat, hitunglah $\int_0^1
\frac{\dd u}{1 + tu} = \frac{\ln(1+t)}{t}$, lalu simpulkan $S_n \sim
\frac{\ln\ln n}{\ln n}$.)}
\end{exercise}

\begin{exercise}[$\star$]\label{exo:b2:comparison:10}
Buktikan kesamaan $(\ln n)^{\ln n} = n^{\ln\ln n}$, lalu urutkan
yang berikut ini menurut $o(\cdot)$ yang menaik di tak hingga, beserta
buktinya: $n^2$, $(\ln n)^{\ln n}$, $2^n$, $n!$, $n^n$.
\end{exercise}

\begin{exercise}[$\star\star$]\label{exo:b2:comparison:11}
(Ekor $\sum 1/k^2$, dua suku) Dengan memakai teleskop eksak
$\sum_{k > n} \frac{1}{k(k+1)} = \frac{1}{n+1}$ dan penguraian
$\frac1{k^2} = \frac{1}{k(k+1)} + \frac{1}{k^2(k+1)}$, buktikan
\[
\sum_{k > n} \frac{1}{k^2}
= \frac1n - \frac{1}{2n^2} + O\Bigl(\frac{1}{n^3}\Bigr).
\]
\end{exercise}

\begin{exercise}[$\star\star\star$]\label{exo:b2:comparison:12}
Misalkan $u_0 = \frac12$ dan $u_{n+1} = u_n + \eu^{-u_n}$. Buktikan bahwa
$u_n \to \infty$, lalu --- dengan menetapkan $v_n = \eu^{u_n}$ dan
menunjukkan $v_{n+1} = v_n + 1 + \frac{1}{2v_n} + O\bigl(v_n^{-2}\bigr)$
--- tegakkan
\[
u_n = \ln n + \frac{\ln n}{2n} + O\Bigl(\frac1n\Bigr).
\]
\end{exercise}

\section{Soal: Bootstrap, dari Euler--Maclaurin ke Bilangan
Prima}

Besaran yang implisit atau terhimpun jarang menyerahkan
asimtotiknya sekaligus; kita menyarikannya lintasan demi lintasan, dan
tiap lintasan mengumpankan taksiran sebelumnya kembali ke relasi
pendefinisinya. Soal akhir pekan ini melatih loop itu pada persamaan yang
segar, membuktikan \emph{rumus Euler--Maclaurin} orde pertama (yaitu
peningkatan trapesium atas perbandingan deret dengan integral, \omterm{def:b2:metric:complete}{lengkap}
dengan batang galat yang cermat), membalik $x\ln x = n$, lalu mencairkan
cek termasyhur metodenya: dari teorema bilangan prima yang diterima tanpa
bukti, diperoleh hukum asimtotik $p_n \sim n\ln n$ bagi bilangan prima ke-$n$.

\begin{problem}[{Soal akhir pekan --- koreksi Euler--Maclaurin
dan asimtotik bilangan prima ke-$n$}]
\label{pb:b2:comparison:1}

\medskip
\noindent\textbf{Bagian I --- Loop bootstrap pada sebuah persamaan
yang segar.}
\begin{enumerate}
  \item Buktikan klaim ketunggalan pada
        \cref{def:b2:comparison:expansion}: jika $f = \sum_{i\leq
        k} c_i\varphi_i + o(\varphi_k) = \sum_{i \leq k}
        c_i'\varphi_i + o(\varphi_k)$ sepanjang skala yang sama, maka
        $c_i = c_i'$ untuk setiap $i$. Lalu doronglah contoh campuran
        di dalam pelajaran satu anak tangga lebih jauh:
        \[
        \frac{1}{x - \ln x} = \frac1x + \frac{\ln x}{x^2} +
        \frac{(\ln x)^2}{x^3} + o\Bigl(\frac{(\ln
        x)^2}{x^3}\Bigr) \qquad (x \to +\infty),
        \]
        lalu jelaskan mengapa tak ada suku $\frac{c}{x^2}$ yang muncul.
  \item Tunjukkan bahwa untuk setiap $n \geq 1$ persamaan $\eu^x + x =
        n$ punya tepat satu penyelesaian real $x_n$, dan bahwa $x_n
        \to +\infty$ dengan $x_n \sim \ln n$.
  \item Bootstraplah dua kali:
        \[
        x_n = \ln n - \frac{\ln n}{n} - \frac{(\ln n)^2}{2n^2} +
        o\Bigl(\frac{(\ln n)^2}{n^2}\Bigr).
        \]
  \item Periksalah secara numerik di $n = 1000$: bandingkan $x_{1000}
        \approx 6.90083$ dengan nilai satu, dua dan tiga suku
        pada pertanyaan 3, sampai lima angka desimal.
\end{enumerate}

\medskip
\noindent\textbf{Bagian II --- Euler--Maclaurin, orde satu.}
\begin{enumerate}[resume]
  \item Buktikan kesamaan kernel trapesium: untuk $g$ berkelas
        $C^2$ pada $\intcc{0}{1}$,
        \[
        \int_0^1 g(t)\,\dd t = \frac{g(0) + g(1)}{2}
        - \frac12\int_0^1 t(1 - t)\,g''(t)\,\dd t
        \]
        \emph{(integralkan $\frac12 t(1-t)g''$ secara parsial dua kali)}.
  \item Misalkan $f$ berkelas $C^2$ pada $\intco{1}{+\infty}$ dengan
        $\int_1^\infty \abs{f''} < \infty$. Tunjukkan bahwa
        \[
        E_n = \sum_{k=1}^{n} f(k) - \int_1^n f -
        \frac{f(1) + f(n)}{2}
        \]
        konvergen ke sebuah konstanta $E$, dengan batas ekornya
        $\abs{E - E_n} \leq \frac18\int_n^\infty\abs{f''}$: itulah
        \emph{rumus Euler--Maclaurin} sampai orde pertama.
  \item Terapkan hal itu pada $f(t) = \frac1t$: buktikan
        \[
        H_n = \ln n + \gamma + \frac{1}{2n} + \varepsilon_n,
        \qquad \abs{\varepsilon_n} \leq \frac{1}{8n^2},
        \]
        yang memperkuat \cref{exo:b2:comparison:3} (kenali
        konstantanya dengan $\gamma$ lewat pembandingan terhadap
        \cref{ex:b2:comparison:harmonic}).
  \item Sarikan koefisien berikutnya: tunjukkan $\varepsilon_n =
        -\frac{1}{12n^2} + o\bigl(\frac1{n^2}\bigr)$
        \emph{(pertambahan $E_n$ adalah
        $\frac12\int_0^1t(1-t)f''(n+t)\dd t =
        \frac1{12}f''(n) + o(f''(n))$; lalu jumlahkan ekornya dengan
        \cref{thm:b2:comparison:seriesintegral})}.
  \item Terapkan pertanyaan 6 pada $f = \ln$: turunkan ulang dalam tiga
        baris kekonvergenan $d_n = \ln n! - (n +
        \frac12)\ln n + n$ (yakni Langkah 1 pada
        \cref{thm:b2:comparison:stirling}), \omterm{def:b2:metric:complete}{lengkap} dengan bonus laju
        galatnya $d_n = d + O\bigl(\frac1n\bigr)$.
  \item Terapkan pertanyaan 6 pada $f(t) = \frac{1}{\sqrt t}$: tunjukkan
        \[
        \sum_{k=1}^{n}\frac1{\sqrt k} = 2\sqrt n + c +
        \frac{1}{2\sqrt n} + O\Bigl(\frac{1}{n^{3/2}}\Bigr)
        \]
        untuk suatu konstanta $c$, lalu nilailah semua sukunya di $n =
        10^4$ (konstantanya $c \approx -1.4604$).
\end{enumerate}

\medskip
\noindent\textbf{Bagian III --- Pembalikan: persamaan $x\ln x =
n$.}
\begin{enumerate}[resume]
  \item Tunjukkan bahwa $x\ln x = n$ punya tepat satu penyelesaian $x_n
        \in \intco{1}{+\infty}$ untuk $n \geq 1$, bahwa $x_n \to
        \infty$, dan bahwa $\ln x_n \sim \ln n$.
  \item Turunkan pembalikan satu sukunya $x_n \sim
        \dfrac{n}{\ln n}$, lalu bootstraplah sekali lagi:
        \[
        \ln x_n = \ln n - \ln\ln n + o(1),
        \qquad
        x_n = \frac{n}{\ln n}\Bigl(1 + \frac{\ln\ln n}{\ln n}
        + o\Bigl(\frac{\ln\ln n}{\ln n}\Bigr)\Bigr).
        \]
  \item Ujilah di $n = 10^6$: akar sejatinya $x \approx
        87\,848$; bandingkan dengan nilai satu suku ($\approx 72\,382$)
        dan dua suku ($\approx 86\,140$), lalu jelaskan
        lambannya keuntungan itu (sebab parameter uraiannya
        $\frac{\ln\ln n}{\ln n}$, yang hanya $\approx 0.19$ di $n = 10^6$).
  \item Kini kita \emph{terima tanpa bukti} teorema bilangan prima:
        banyaknya bilangan prima $\pi(x)$ yang $\leq x$ memenuhi $\pi(x)
        \sim \frac{x}{\ln x}$ ketika $x \to \infty$ (yang dibuktikan
        secara jujur pada jilid Tahun ke-3). Dengan menulis $p_n$ untuk
        bilangan prima ke-$n$, berilah alasan bahwa $\pi(p_n) = n$, lalu
        jalankan pembalikan pada pertanyaan 11--12 untuk membuktikan
        \[
        p_n \sim n \ln n .
        \]
  \item Panennya: (a) tunjukkan $\sum_{k \leq n} p_k \sim
        \frac{n^2\ln n}{2}$ \emph{(bandingkan $\sum k\ln k$ dengan
        $\int t\ln t\,\dd t$)}; (b) hitunglah kira-kira berapa
        peluang sebuah bilangan bulat acak seragam berangka $100$
        merupakan bilangan prima ($\ln 10^{100} \approx 230.26$: kira-kira
        satu di antara $230$).
\end{enumerate}

\medskip
\noindent\textbf{Bagian IV --- Metodenya diekspor: $x\tan x =
1$.}
\begin{enumerate}[resume]
  \item Tunjukkan bahwa untuk tiap $n \geq 1$ persamaan $\tan x =
        \frac1x$ punya tepat satu penyelesaian $x_n$ di
        $\intoo{n\pi}{\,n\pi + \frac\pi2}$, dan bahwa $z_n = x_n
        - n\pi \to 0^+$.
  \item Satu suku: $z_n \sim \dfrac{1}{n\pi}$.
  \item Tunjukkan bahwa uraian $z_n$ \emph{tidak} punya
        suku $\frac{c}{n^2}$: yakni $z_n = \frac1{n\pi} +
        O\bigl(\frac{1}{n^3}\bigr)$.
  \item Tiga suku: dengan memakai $\arctan u = u - \frac{u^3}3 +
        O(u^5)$ dan $\frac1{x_n} = \frac{1}{n\pi} -
        \frac{z_n}{(n\pi)^2} + O(n^{-3}\cdot z_n^2)$, buktikan
        \[
        x_n = n\pi + \frac{1}{n\pi} -
        \frac{4}{3\pi^3 n^3} + o\Bigl(\frac{1}{n^3}\Bigr).
        \]
  \item Periksalah di $n = 3$: akar sejatinya $x_3 \approx 9.5293344$;
        bandingkan nilai satu suku dengan nilai tiga sukunya, lalu
        bandingkan dalam satu kalimat dengan $\tan x = x$ di dalam
        pelajaran (\cref{ex:b2:comparison:tan}): di mana tiap barisannya
        duduk pada jendelanya, dan mengapa.
\end{enumerate}

\medskip
\noindent\textbf{Bagian V --- Bootstrap dinamis, aturan
mainnya, dan rangkuman.}
\begin{enumerate}[resume]
  \item Misalkan $u_0 \in \intoo{0}{\pi}$ dan $u_{n+1} = \sin u_n$.
        Tunjukkan $u_n \to 0$ secara menurun, lalu hitunglah limit
        $\dfrac{1}{u_{n+1}^2} - \dfrac{1}{u_n^2}$ \emph{(uraikan
        $\sin^{-2}$ lewat $\sin u = u - \frac{u^3}6 + o(u^3)$)}.
  \item Turunkan, lewat rata-rata Cesàro (jilid Tahun ke-1), hasil klasik
        \[
        u_n \sim \sqrt{\frac{3}{n}} .
        \]
  \item (Numerik yang tersahkan) Dengan memakai batas cermat pada
        pertanyaan 7, tunjukkan bahwa menilai $\ln n + \gamma +
        \frac1{2n}$ di $n = 10^6$ menghasilkan $H_{10^6}$ dengan galat
        paling banyak $1.25\cdot10^{-13}$ --- yakni jumlah sejuta suku
        yang terhitung sampai tiga belas angka oleh tiga suku saja.
  \item (Aturan mainnya) Buktikan atau sangkallah, beserta buktinya atau
        contoh penyangkalnya: (a) jika $u_n \sim v_n \to +\infty$ maka
        $\ln u_n \sim \ln v_n$; (b) jika $u_n \sim v_n$ maka
        $\eu^{u_n} \sim \eu^{v_n}$; (c) jika $f \sim g$ di
        $+\infty$ (dengan $f, g$ dapat diturunkan) maka $f' \sim g'$.
  \item (Rangkuman) Satu kalimat untuk masing-masing: loop bootstrap pada
        \cref{met:b2:comparison:implicit} sebagaimana dipakai pada Bagian
        I, III, dan IV; apa yang ditambahkan koreksi trapesium kepada
        \cref{thm:b2:comparison:seriesintegral}; mengapa pembalikan
        $x\ln x$ persis menjadi jembatan dari $\pi(x)$ ke
        $p_n$; dan aturan mana pada pertanyaan 24 yang melindungi langkah
        yang mana. Sebutkan kedua puncaknya: rumus Euler--Maclaurin
        (orde pertama), dan hukum asimtotik bilangan prima ke-$n$.

\end{enumerate}
\end{problem}
